% !PRINTED 570
\weberpage{570}

%<<CONT
\weq{8}{P_\mu(x) = \Pi\left[x - \wp\left(\frac{\nu}{\mu}\right)\right],}

and $P_\mu(x)$ is of degree $m - 1$.

In the product $P_\mu(x)$ every factor $x - g$ occurs twice, except when $g$ is equal to one of the $e$, that is, when the congruence (7) can be satisfied for a $\nu$ different from zero.

1) If $\mu$ is relatively prime to 2, and hence if $m$ is \gesperrt{odd}, then the congruence (7) has only the root $\nu = 0$. In that case $P_\mu(x)$ is a square, and if $S(x)$ is a function of degree $\frac{1}{2}(m - 1)$,

\weq{9}{P_\mu(x) = S^2, \quad m \equiv 1 \ (\mathrm{mod}\ 2).}

If $\mu$ has a divisor in common with 2, three cases are to be distinguished.

2) If $\mu$ is divisible by 2, then $\nu$, if it is to satisfy the congruence (7), must be a multiple of $\frac{1}{2}\mu$; and since there are in $\Omega$ three numbers different from zero that are incongruent with respect to the modulus 2, (7) has three roots. Accordingly the three values $e_1, e_2, e_3$ occur among the $g$, and we have:

\weq{10}{P_\mu(x) = \wp'(u)^2 S^2, \quad m \equiv 0 \ (\mathrm{mod}\ 2),}

where $S$ is an integral function of degree $\frac{1}{2}m - 2$.

3) If 2 splits in the field $\Omega$ into two (equal or different) prime ideals, that is, if $\Delta \equiv 0 \ (\mathrm{mod}\ 4)$ or $\equiv 1 \ (\mathrm{mod}\ 8)$, then $\mu$ can be divisible by one of these prime factors without being divisible by 2. In that case the congruence (7) has only one root different from zero, and among the numbers $g$ only one of the three numbers (6), say $e$, occurs. In this case

\weq{11}{P_u(x) = [\wp(u) - e]\, S^2,} % [sic] The source prints P_u; the surrounding formulas and prose require P_\mu.

where $S$ is a function of degree $\frac{1}{2}m - 1$.

With regard to the two special cases $g_3 = 0$ and $g_2 = 0$ the following is still to be remarked.

% [sic] The source sets j(omega)=0 here. By (3), § 150, g_3=0 gives j=27*64; the values in cases 4) and 5) are interchanged and are retained as printed.
4) If $g_3 = 0$, then $\Delta = -4$, $j(\omega) = 0$, and from the differential equation §\,46, (18) there follows:

\[
\wp'(u)^2 = 4\wp(u)^3 - g_2 \wp(u),
\]

by which, having regard to the initial term of the development, the $\wp$-function is uniquely determined.
